\documentclass[11pt,a4paper]{article}
\usepackage[T1]{fontenc}
\usepackage[margin=28mm]{geometry}
\usepackage{amsmath,amssymb,amsthm,booktabs,longtable}
\usepackage{microtype}
\usepackage[hidelinks]{hyperref}
\usepackage{url}
\newtheorem{theorem}{Theorem}[section]
\newtheorem{lemma}[theorem]{Lemma}
\newtheorem{proposition}[theorem]{Proposition}
\newtheorem{corollary}[theorem]{Corollary}
\theoremstyle{definition}
\newtheorem{definition}[theorem]{Definition}
\newcommand{\F}{\mathbb F_5}
\newcommand{\B}{\mathcal B}
\newcommand{\HH}{\mathcal H}
\newcommand{\RR}{\mathcal R}
\newcommand{\cl}{\operatorname{cl}}
\newcommand{\si}{\operatorname{si}}
\newcommand{\PG}{\operatorname{PG}}
\newcommand{\Aut}{\operatorname{Aut}}
\setlength{\emergencystretch}{2em}
\setcounter{tocdepth}{1}
\title{The rank-four, eleven-element excluded minors\protect\\
for representability over $\mathbb F_5$}
\author{}
\date{A computational classification\\Revised exposition, 5 November 2026}

\begin{document}
\maketitle

\begin{abstract}
There are exactly 35 isomorphism classes of rank-four, eleven-element
excluded minors for representability over $\mathbb F_5$. The proof is an
exhaustive construction. Every such matroid has a representable
ten-element deletion. For each possible deletion, we recover all
nonrepresentable extensions whose other deletions are representable,
then test their contractions. The essential observation is that, when
the ground set has more than twice the rank many elements, compatible
basis data for the single-element deletions already determine a
matroid. This turns the extension search into an exact intersection of
finite tables. We prove that each stage preserves every possible
excluded minor, explain how the tables are constructed and exhausted,
and give the computational results that complete the classification.
Separate certificates verify nonrepresentability and all immediate
minors of each of the 35 resulting matroids.
\end{abstract}

\tableofcontents
\clearpage

\section{The theorem and the structure of its proof}\label{sec:classification}

A matroid is \emph{quinary} if it is representable over
$\F=\mathbb Z/5\mathbb Z$. An \emph{excluded minor} for quinary matroids
is a nonquinary matroid every proper minor of which is quinary. We
classify these obstructions at one specified rank and size.

\begin{theorem}[Computational classification]\label{thm:main}
There are exactly 35 isomorphism classes of rank-four, eleven-element
excluded minors for representability over $\F$. Representatives
$M_1,\ldots,M_{35}$ are specified in Appendix~\ref{sec:catalogue}.
\end{theorem}

The defining condition has a useful finite form. A matroid $M$ is an
excluded minor precisely when
\[
 M\text{ is nonquinary},\qquad
 M\setminus a\text{ and }M/a\text{ are quinary for every }a\in E(M).
\]
Indeed, every proper minor is a minor of one of these immediate
deletions or contractions, and field representability is closed under
taking minors. For an eleven-element matroid there are therefore 22
immediate minors to consider. This criterion is easy to apply to a
given matroid. The main difficulty is to find a finite list that is
guaranteed to contain \emph{every} matroid satisfying it.

We construct such a list by reversing deletion. Choose an element
$e$ of an unknown excluded minor $M$ and put $N=M\setminus e$. The
matroid $N$ is quinary. We call it a \emph{parent}, since the task is
to recover $M$ by adding $e$ to it. The proof proceeds in three steps.
\begin{enumerate}
\item \emph{Find all possible parents.} We prove that $N$ must be a
connected simple rank-four matroid on ten elements, and generate a
complete quinary catalogue of these matroids
(Section~\ref{sec:parents}).
\item \emph{Recover all possible extensions of each parent.} For a
fixed $N$, we construct exactly the nonquinary extensions $M$ for which
every deletion $M\setminus a$ is quinary. The recovery theorem is
proved in Section~\ref{sec:local}; its finite tables and their exhaustive
solution are developed in Section~\ref{sec:tables}.
\item \emph{Impose the remaining minor conditions.} We test every
contraction of every recovered matroid and identify isomorphic
outputs. This leaves precisely the desired excluded minors
(Section~\ref{sec:finish}).
\end{enumerate}

The second step is the heart of the argument. A quinary extension of
$N$ can be obtained by appending a column to a representation of $N$;
the desired $M$, however, is nonquinary. Instead we construct its
\emph{different deletions} from representations, and require their
basis data to agree wherever they overlap. A local-to-global lemma
then shows that these compatible data form a matroid. We can thus
enumerate nonrepresentable matroids using only finite lists of
representable ones.

Two inclusions will finish the proof. The construction must contain
every excluded minor, giving an upper bound on their number. Each
listed object must also be an excluded minor, and the objects must be
pairwise nonisomorphic, giving the matching lower bound. The
computation in Section~\ref{sec:computation} supplies both: the
exhaustive construction leaves 35 classes, and separate certificates
verify all 35. Its numerical results are organized there according
to the three steps above; implementation records are collected in
Appendix~\ref{sec:reproduction}.

For comparison, Brettell's classification through ten elements gives
5, 92, 219, and 130 rank-four excluded minors at sizes seven through
ten \cite{brettell}. We reproduce those smaller catalogues as checks
on the method. The use of compatible extensions and of inequivalent
representations has precedents in the computations of Brettell and
Pendavingh \cite{brettell-pendavingh}. Here the proof of coverage rests
on generating all possible parents and imposing compatibility with
every deletion.

\section{Step 1: finding every possible parent}\label{sec:parents}

This step establishes the starting point of the search: a complete
list of matroids $N$ that could occur as $M\setminus e$. We first
determine which parents are needed, then prove that the represented
generation procedure reaches all of them.

\subsection{Why connected simple parents suffice}

We use the standard notation $r_M(X)$ for the rank of a subset and
$\B(M)$ for the bases of $M$. All matroids are finite. A simple
matroid has no loops or parallel pairs.

\begin{lemma}[Possible deletions]\label{lem:parents}
If $M$ is a rank-four, eleven-element excluded minor for quinary
matroids, then for every $e\in E(M)$ the deletion $M\setminus e$ is
a connected simple rank-four quinary matroid on ten elements.
\end{lemma}
\begin{proof}
An excluded minor for representability over a field is 3-connected
when it has at least four elements. To recall the reason,
representability over a fixed field is closed under direct sums and
2-sums. A disconnected matroid decomposes as a direct sum; a connected
matroid with a 2-separation decomposes as a 2-sum. The summands are
isomorphic to proper minors \cite{oxley}. If the matroid were an
excluded minor, all summands would be representable, forcing the
matroid itself to be representable. For the 2-sum construction, choose
representations in which the common element is a unit vector,
identify its one-dimensional span in the two vector spaces, and
delete its column.

Thus $M$ is 3-connected. It has no loops, parallel pairs, or coloops,
so $N=M\setminus e$ is simple and has rank four. It is quinary by
minor-minimality. It remains to prove that $N$ is connected. If
$E(N)=A\mathbin{\dot\cup}B$ were a nontrivial direct-sum partition,
choose the names so that $|A|\geq2$. Since $r(N)=r(M)$, we would have
\[
 r_M(A)+r_M(B)=r(M),\qquad
 r_M(A)+r_M(B\cup\{e\})-r(M)\leq1.
\]
Both $A$ and $B\cup\{e\}$ have at least two elements. This contradicts
3-connectivity of $M$.
\end{proof}

Let $P=\{0,\ldots,9\}$ and take a new label $e=10$. Let
$\mathcal P$ contain one labelled representative on $P$ of each
connected simple rank-four quinary matroid on ten elements. For each
$N\in\mathcal P$, define $\mathcal D(N)$ to be the set of simple
rank-four matroids $M$ on $P\cup\{e\}$ satisfying
\[
 \begin{gathered}
 M\setminus e=N,\qquad M\text{ is nonquinary},\\
 M\setminus a\text{ is quinary for every }a\in P\cup\{e\}.
 \end{gathered}
\]
We call these \emph{deletion-minimal nonquinary extensions} of $N$.
The word ``deletion'' is essential: contractions have not yet been
tested.

\begin{proposition}[Reduction of the classification]\label{prop:reduction}
The rank-four, eleven-element quinary excluded minors are exactly,
up to isomorphism, the members of
\[
 \bigcup_{N\in\mathcal P}
 \{M\in\mathcal D(N):M/a\text{ is quinary for every }a\in E(M)\}.
\]
\end{proposition}
\begin{proof}
Every matroid in the displayed set is nonquinary and has quinary
immediate deletions and contractions, so it is an excluded minor.
Conversely, let $M$ be any excluded minor of the specified rank and
size. By Lemma~\ref{lem:parents}, a deletion $M\setminus e$ is
isomorphic to a member $N$ of $\mathcal P$. Relabel this deletion as
$N$ and its missing element as $e$. The resulting matroid belongs
to $\mathcal D(N)$, and all its contractions are quinary. Its
isomorphism class therefore occurs in the displayed set.
\end{proof}

This proposition specifies exactly what the rest of the computation
must produce. It suffices to construct $\mathcal P$, compute each
$\mathcal D(N)$ completely, and apply the contraction test. In
particular, a parent may be discarded for being disconnected, but
not merely for failing to be 3-connected.

\subsection{Generating a complete represented catalogue}

A full-rank $r\times k$ matrix over $\F$ represents a matroid whose
bases are its nonsingular $r$-column submatrices. In a simple
representation, its columns specify distinct points of
$\PG(r-1,5)$ and span the ambient space. We represent a projective
point by the vector whose first nonzero entry is 1. There are
$(5^r-1)/4$ such points, in particular 156 in rank four.

Starting with $I_r$, append every projective point not already used,
and repeat until the required size is reached. Isomorphic repetitions
can be removed during growth, provided the equivalence used respects
the \emph{representation}. Two matrices are \emph{field-isomorphic}
when one can be obtained from the other by invertible row operations,
nonzero column scalings, and a permutation of columns. Abstractly
isomorphic matroids can have representations that are not
field-isomorphic. All such represented classes must be retained.

\begin{proposition}[Completeness of represented generation]\label{prop:generation}
The procedure just described, retaining one matrix from each
field-isomorphism class at each size, generates every simple
rank-$r$ quinary representation up to field isomorphism.
\end{proposition}
\begin{proof}
Take any simple rank-$r$ representation. Choose $r$ columns forming
a basis, transform them to $I_r$, and order the remaining columns.
Appending those columns in that order gives a chain of simple
rank-preserving represented extensions from $I_r$ to the chosen
matrix. Inductively, each intermediate representation is
field-isomorphic to a retained one. Such an isomorphism carries all
possible appended projective points bijectively to possible appended
points of the retained representation, and extends the permutation
of old labels by fixing the new label. Thus replacing an intermediate
matrix by its retained representative cannot lose any later
represented extension class. Induction on the number of columns
proves the assertion.
\end{proof}

After a represented level is complete, forgetting the matrices and
deduplicating by abstract matroid isomorphism gives the required
parent classes. We generate all simple rank-four represented levels
through size ten in this way, then retain the connected ten-element
parents for $\mathcal P$. The same procedure in rank three provides
the catalogue that will be needed to test contractions in
Section~\ref{sec:finish}.

Step 1 is now justified: every possible deletion of an excluded minor
occurs in $\mathcal P$. We next explain how to recover its possible
nonquinary extensions. This requires more than appending columns to
the parent itself; it requires matching the representable deletions
of the prospective extension.

\section{Step 2: recovering an extension from its deletions}\label{sec:local}

Fix one parent $N\in\mathcal P$. Our goal in this section is an exact
description of $\mathcal D(N)$ that can be computed from finite lists
of quinary matroids. We first describe the unknown data of an
extension, then express the conditions on its deletions in those
data. The main proposition will show that these conditions already
enforce the matroid axioms.

\subsection{The unknown basis data}

The deletion equation $M\setminus e=N$ fixes the status of every
four-set avoiding $e$: it is a basis of $M$ exactly when it is a
basis of $N$. The only unknown four-sets are therefore
\[
 T\cup\{e\},\qquad T\in\binom P3.
\]
There are $\binom{10}{3}=120$ of these, compared with the
$\binom{10}{4}=210$ already fixed four-sets. Introduce a Boolean
variable $x_T$ for each triple, with $x_T=1$ meaning that
$T\cup\{e\}$ is a basis. An assignment $x$ specifies a nonempty
four-subset family
\[
 \B_N(x)=\B(N)\ \cup
          \{T\cup\{e\}:T\in\binom P3,\ x_T=1\}.
\]
For the moment this is only a proposed basis family; an arbitrary
assignment need not define a matroid.

Consider how another deletion reads these variables. For $f\in P$,
the four-sets of $M\setminus f$ that contain $e$ are indexed by the
triples avoiding $f$. Write
\[
 \pi_f(x)=x\big|_{\binom{P\setminus\{f\}}3}
\]
for the restriction to those coordinates. The four-sets avoiding
both $e$ and $f$ are already determined by $N\setminus f$. Thus
$\pi_f(x)$ is exactly the extra basis data needed to specify
$M\setminus f$ as an extension of $N\setminus f$.

For example, the bit $x_{\{0,1,2\}}$ records whether
$\{0,1,2,10\}$ is a basis. This bit is visible in the deletion of
3 and in the deletion of 4, but disappears in the deletion of 0.
If independently chosen extensions for the deletions of 3 and 4
disagree on that bit, they cannot be deletions of one matroid. A
single assignment $x$ enforces agreement on every such shared
four-set at once.

To express which deletions are permitted, let $Q$ be any fixed
labelled simple rank-four quinary matroid on a set $U$. Define its
\emph{represented extension table} $\RR(Q)$ to be the set of Boolean
patterns $z$ on $\binom U3$ for which
\[
 \B(Q)\ \cup\{T\cup\{e\}:z_T=1\}
\]
is the basis family of a simple rank-four quinary extension of $Q$.
All labels in this definition are fixed. In
Section~\ref{sec:tables} we will construct this table completely
from representations of $Q$.

The meaning of table membership is now precise:
\begin{itemize}
\item $\pi_f(x)\in\RR(N\setminus f)$ says that the proposed deletion
avoiding $f$ is simple, rank four, and quinary;
\item $x\in\RR(N)$ says that the proposed full extension of $N$ is
itself simple, rank four, and quinary.
\end{itemize}
The first statement is meaningful because a connected parent has no
coloops, so every $N\setminus f$ still has rank four.

Consequently the assignments we want should form the finite set
\[
 \mathcal S(N)=
 \bigl\{x\in\{0,1\}^{\binom P3}:
       \pi_f(x)\in\RR(N\setminus f)\text{ for every }f\in P\bigr\}
       \setminus\RR(N).
\]
The intersection requires all the other deletions to be quinary;
the subtraction removes the extensions that are quinary themselves.
What still needs proof is that an assignment satisfying these
conditions actually defines a matroid. Without that implication,
the displayed expression would merely describe consistent tables
of four-subsets.

\subsection{Why compatible deletions force the matroid axioms}

The basis-exchange axiom involves two bases and hence at most twice
the rank many elements. If the ground set is larger, the whole
exchange question can be seen inside one deletion. This gives the
needed local-to-global principle.

\begin{lemma}[Deletion compatibility]\label{lem:deletion}
Let $|E|=n>2r$ and let $\varnothing\ne\B\subseteq\binom Er$.
Suppose that for each $a\in E$ the family of members of $\B$ avoiding
$a$ is the basis family of a rank-$r$ matroid on $E\setminus\{a\}$.
Then $\B$ is the basis family of a rank-$r$ matroid on $E$.
\end{lemma}
\begin{proof}
Take $A,B\in\B$ and $b\in A\setminus B$. We must find
$c\in B\setminus A$ such that
$(A\setminus\{b\})\cup\{c\}\in\B$. Since
$|A\cup B|\leq2r<n$, there is an element $a\notin A\cup B$.
Both $A$ and $B$ are bases of the specified deletion avoiding $a$.
Basis exchange in that matroid supplies the required $c$, and its
exchanged basis also belongs to $\B$. This proves every ordered
instance of basis exchange. Together with nonemptiness and the
common size of the members, it proves that $\B$ is a matroid basis
family.
\end{proof}

Here $n=11$ and $r=4$, so the lemma applies. In particular, the
representations witnessing different deletions do not have to fit
together as matrices. They need only prescribe compatible bases.
The lemma then supplies the abstract matroid, even when it has no
representation over $\F$.

\begin{proposition}[Exact recovery of deletion-minimal extensions]\label{prop:local}
For a fixed parent $N\in\mathcal P$, the map
$x\mapsto M_x$, where $\B(M_x)=\B_N(x)$, is a bijection from
$\mathcal S(N)$ to $\mathcal D(N)$.
\end{proposition}
\begin{proof}
First take $x\in\mathcal S(N)$. The proposed basis family
$\B_N(x)$ is nonempty because it contains $\B(N)$. Its deletion
avoiding $e$ is the rank-four quinary matroid $N$. For each $f\in P$,
the condition $\pi_f(x)\in\RR(N\setminus f)$ says that its deletion
avoiding $f$ is also the basis family of a simple rank-four quinary
matroid. Lemma~\ref{lem:deletion} therefore makes $\B_N(x)$ a
matroid basis family.

Call this matroid $M_x$. Each pair of elements lies in one of its
simple deletions, so $M_x$ is simple. All its deletions are quinary
by construction. If $M_x$ were quinary, its pattern would belong
to $\RR(N)$ by the definition of that table, contrary to the
subtraction in $\mathcal S(N)$. Hence $M_x\in\mathcal D(N)$.

Conversely, take $M\in\mathcal D(N)$. Its bases containing $e$
give a unique pattern $x$. For every $f\in P$, the matroid
$M\setminus f$ is a simple quinary extension of $N\setminus f$;
it has rank four since it contains the rank-four restriction
$N\setminus f$. Thus $\pi_f(x)\in\RR(N\setminus f)$. Since
$M$ is nonquinary, $x\notin\RR(N)$. Therefore
$x\in\mathcal S(N)$, and $M=M_x$. Finally, different patterns
give different basis families, proving injectivity.
\end{proof}

This proposition is the bridge between the mathematical problem
and the computation. It accounts for all four requirements on an
output: validity as a matroid, simplicity, representability of every
deletion, and nonrepresentability of the full matroid. To finish
Step 2, it remains to build the tables $\RR(Q)$ without omissions
and exhaust the finite conditions defining $\mathcal S(N)$.

\section{Constructing and exhausting the extension tables}\label{sec:tables}

Proposition~\ref{prop:local} reduces each parent search to a finite
problem about quinary extensions of nine- and ten-element matroids.
There are two obligations: every represented extension must occur
in its table, and every simultaneous solution of the table
conditions must be considered. We address them in that order.

\subsection{Putting all representations on the required labels}

The represented catalogues of Section~\ref{sec:parents} retain all
field-isomorphism classes. A table $\RR(Q)$, however, describes
extensions of one \emph{labelled} matroid $Q$. To use the catalogue
for this purpose, we need all representations on those labels, up
to row operations and column scaling alone. We call the latter
equivalence \emph{labelled projective equivalence}.

First choose one matroid isomorphism aligning each unlabelled
represented class of $Q$ with the chosen labels. There may be many
alignments, and they can give different labelled extension patterns.
The missing alignments are recovered by applying all automorphisms
of $Q$ to the aligned matrices.

\begin{lemma}[Complete labelled representations]\label{lem:labelled}
Align one matrix from every field-isomorphism class representing
$Q$ with the fixed labels of $Q$, and close the resulting set under
$\Aut(Q)$. Up to labelled projective equivalence, this gives every
representation of the labelled matroid $Q$.
\end{lemma}
\begin{proof}
Take any representation on those labels and forget its labels.
It belongs to one of the represented classes in the complete
catalogue. Choose a field isomorphism from the catalogue matrix
to this representation. It induces an alignment of the underlying
matroid with $Q$. The chosen alignment of that same catalogue
matrix differs from it by an automorphism of $Q$. Applying that
automorphism to the aligned matrix therefore reaches the desired
representation, up to row operations and column scaling.
\end{proof}

The computation closes under generators of the full automorphism
group, which reaches the same orbit as applying the entire group.
Repeated labelled projective classes are recognized by a canonical
matrix normalization. The normalization and its proof are given in
Appendix~\ref{sec:normalization}; its purpose is to remove repetitions
without discarding a representation class. Denote the resulting
complete labelled set of matrices by $\mathcal L(Q)$.

\subsection{Building a table by appending projective points}

For each $V\in\mathcal L(Q)$ append every point $v$ of $\PG(3,5)$
whose direction is not already a column of $V$. The resulting
matrix has rank four, and its column matroid is simple. Record the
pattern
\[
 \rho(V,v)_T=1\quad\Longleftrightarrow\quad
             \det[\,V_T\mid v\,]\ne0,
             \qquad T\in\binom{E(Q)}3,
\]
where $V_T$ consists of the three old columns indexed by $T$.
Collect all these rows, merging duplicates.

\begin{proposition}[Completeness of the extension table]\label{prop:tables}
The rows obtained in this way are exactly $\RR(Q)$.
\end{proposition}
\begin{proof}
Every appended point gives a simple rank-four quinary extension,
so every constructed row belongs to $\RR(Q)$. Conversely, take
any extension represented by a row of $\RR(Q)$, and choose a
matrix representing it. Its restriction to $Q$ is labelled
projectively equivalent to a member $V$ of $\mathcal L(Q)$, by
Lemma~\ref{lem:labelled}. Apply that equivalence to the entire
matrix. The new column is a nonzero projective point not already
used, since the extension is simple. That point is among those
appended to $V$, so the extension's row occurs in the construction.
\end{proof}

This proves the first obligation of the section. In particular, a
deletion $N\setminus f$ is matched to its catalogue representative
and all its matrices are aligned back to the actual labels
$P\setminus\{f\}$ before its table is formed. Equality of bits in
different tables consequently refers to the same elements, as
required by Proposition~\ref{prop:local}.

All determinant calculations take place in $\F$. For a fixed old
triple the determinant is a linear form in $v$, so its cofactors
allow all 156 projective points to be tested in one finite
calculation. This speeds up table construction while computing
exactly the determinant bits in the displayed definition.

\subsection{Compressing the coordinates without changing the solutions}

The 120 triple coordinates give the simplest proof of the recovery
theorem, but many must agree. We now identify those repetitions to
obtain the smaller problems actually solved by the program.

Let $M$ be any rank-preserving matroid extension of $N$. If a triple
$T$ is dependent in $N$, then $r_M(T\cup\{e\})\leq3$, so its bit
is zero. If $T$ is independent and $H=\cl_N(T)$, then $H$ is a
rank-three hyperplane and
\[
 \cl_M(T)=\cl_M(H).
\]
Indeed, $T\subseteq H\subseteq\cl_M(T)$, which proves both closure
inclusions. Hence
\[
 T\cup\{e\}\in\B(M)
 \quad\Longleftrightarrow\quad e\notin\cl_M(H).
\]
Thus independent triples spanning the same hyperplane have the
same bit. Since every hyperplane has a spanning triple, one bit
$y_H$ for each rank-three hyperplane $H$ determines all the triple
bits of a genuine extension.

Write $\HH(N)$ for these hyperplanes. In a table for the deletion
of $f$, the visible hyperplanes are
\[
 C_f=\{\cl_N(T):T\in\binom{P\setminus\{f\}}3,
                              \ r_N(T)=3\}.
\]
The closure of such a triple in $N\setminus f$ is
$\cl_N(T)\setminus\{f\}$. Consequently every row of
$\RR(N\setminus f)$ is constant on triples assigned the same
coordinate in $C_f$, and is zero on dependent triples. Projecting
onto one bit per visible hyperplane therefore loses no row
information. Let $A_f\subseteq\{0,1\}^{C_f}$ be this projected
table, and let $A\subseteq\{0,1\}^{\HH(N)}$ be the analogous
projection of $\RR(N)$.

The problem actually solved is
\[
 \widehat{\mathcal S}(N)=
 \{y\in\{0,1\}^{\HH(N)}:
                  y|_{C_f}\in A_f\text{ for every }f\in P\}
                  \setminus A.
\]
Expanding $y_H$ back to every independent triple spanning $H$,
and inserting zeros on dependent triples, takes each solution to
one in $\mathcal S(N)$. Conversely,
Proposition~\ref{prop:local} makes every solution of
$\mathcal S(N)$ a genuine matroid extension, so the closure
argument above shows that it has a unique such compression.
Thus $\widehat{\mathcal S}(N)$ and $\mathcal S(N)$ have exactly the
same solutions in different coordinates.

Every hyperplane coordinate is seen by at least one deletion:
choose $f\notin H$ and a spanning triple of $H$. That triple avoids
$f$, so $H\in C_f$. The union of the coordinate sets $C_f$ is
therefore all of $\HH(N)$.

\subsection{Why the finite join is exhaustive}

It remains to impose all the tables simultaneously. For relations
$U$ and $V$ on coordinate sets $S$ and $T$, their \emph{natural join}
is the finite relation
\[
 U\Join V=\{u\cup v:u\in U,\ v\in V,
                            \ u|_{S\cap T}=v|_{S\cap T}\}.
\]
In words, take every pair of rows that agree on their shared
coordinates and combine them into a row on $S\cup T$. Grouping
rows by their overlap values implements exactly this operation.

Start with the single empty assignment and join the ten relations
$A_f$, one at a time. After processing any collection of deletions,
the current rows are exactly the assignments on the coordinates
seen so far that satisfy all the processed relations. This is
true initially, and the definition of join proves the induction
step. After all ten joins, every coordinate has been seen, so
the rows are exactly
\[
 \{y:y|_{C_f}\in A_f\text{ for every }f\in P\}.
\]
Removing the rows in $A$ now gives
$\widehat{\mathcal S}(N)$ and hence, by the recovery theorem,
every member of $\mathcal D(N)$.

The implementation may change the join order and remove rows of
$A$ as soon as all coordinates have appeared. Neither change
alters the result: a completed assignment cannot be changed by
later joins, only accepted or rejected. Resource limits are
handled separately. An interrupted or bounded calculation is
recorded as \texttt{UNKNOWN}, and cannot count as an exhaustive
parent search. Only a finished calculation establishes that
every solution has been considered. A second method solves the
same finite relations using SAT; its encoding is described in
Appendix~\ref{sec:sat}.

Step 2 is therefore justified: complete representation lists give
complete tables, and completed joins give exactly
$\mathcal D(N)$. At this stage every output is nonquinary and all
its deletions are quinary. The sole remaining minor conditions
are the contractions.

\section{Step 3: testing contractions and certifying the list}\label{sec:finish}

Proposition~\ref{prop:reduction} tells us how to finish the
classification once the sets $\mathcal D(N)$ have been recovered.
We must test every contraction and then pass from labelled
matroids to isomorphism classes. We also verify each final object
by certificates that can be checked without rerunning the search.

\subsection{The remaining minor conditions}

For a simple rank-four matroid $M$, every contraction $M/a$ has
rank three. It may have parallel elements. Its
\emph{simplification} $\si(M/a)$ deletes loops and keeps one
element of each parallel class.

\begin{lemma}[The contraction test]\label{lem:contractions}
For $M\in\mathcal D(N)$, the matroid $M$ is an excluded minor if
and only if every $\si(M/a)$ is a quinary rank-three matroid.
\end{lemma}
\begin{proof}
A matroid is quinary if and only if its simplification is. One
direction follows because a simplification is a restriction; for
the other, extend a representation of the simplification by zero
columns for loops and scalar copies of columns for parallel
elements. Thus the stated condition is equivalent to all
contractions of $M$ being quinary. Membership in $\mathcal D(N)$
already supplies nonrepresentability of $M$ and representability
of all its deletions. The immediate-minor criterion from
Section~\ref{sec:classification} now gives the result.
\end{proof}

This is where the rank-three represented catalogue from Step 1
is used. Each simplification has at most ten elements. A simple
rank-three matroid on three elements is free and is represented
by $I_3$. At sizes four through ten, exact isomorphism comparison
with the corresponding complete quinary catalogue decides the
test. In particular, this step needs no assumption about a
classification of rank-three excluded minors of arbitrary size.

After applying this test, Proposition~\ref{prop:reduction} shows
that the retained matroids are exactly the required excluded
minors, with possible repetitions from their different parents
or labellings. Exact matroid isomorphism removes these repetitions.
The final distinctness check also uses colored incidence graphs:
one color consists of element vertices, the other of vertices
for all bases or, when fewer, all nonbases, with edges given by
membership. A color-preserving graph isomorphism induces exactly
a matroid isomorphism, since either subset family determines the
bases. Canonical forms of these graphs therefore give an exact
second check of pairwise nonisomorphism.

\subsection{Certifying that every listed object is an excluded minor}

The construction proves coverage of the target class. For the
opposite inclusion, each final basis family is validated and
given two kinds of evidence: representations of all its immediate
minors, and a proof that the matroid itself has no representation.

The positive evidence is direct. For each $a\in E(M)$ we save
matrices for $M\setminus a$ and $M/a$, in the order of their
remaining labels. An independent modular Gaussian-elimination
routine computes the basis family of each matrix and compares
it with that of the required minor. It checks every subset of
the appropriate size, so agreement establishes equality of the
represented matroid and the specified minor. Sage performs a
second matrix check. These witnesses certify all the proper-minor
conditions using only finite-field arithmetic.

Nonrepresentability requires a different form of evidence. We
turn existence of a representation into a finite propositional
formula $\Phi_M$ and check a certificate that this formula is
unsatisfiable. Here is the mathematical condition encoded by
the formula. Choose a basis $B=(b_1,\ldots,b_r)$ of $M$ and
reduce its columns to the identity. The remaining entries form
a block $Z$, and fundamental exchanges prescribe its support:
\[
 Z_{ij}\ne0\quad\Longleftrightarrow\quad
 (B\setminus\{b_i\})\cup\{j\}\in\B(M),\qquad j\notin B.
\]
The unknown nonzero entries range over $\{1,2,3,4\}$. For every
$r$-subset $S$, the determinant of the columns on $S$ must be
nonzero exactly when $S\in\B(M)$. These conditions impose the
entire basis family.

\begin{proposition}[Meaning of the nonrepresentability certificate]\label{prop:certificate}
The exact finite-field encoding $\Phi_M$ of these determinant
conditions is satisfiable if and only if $M$ is quinary.
\end{proposition}
\begin{proof}
Any representation of $M$ can be reduced on $B$ to identity form.
Its remaining entries have the prescribed support and satisfy
every determinant condition, so they give a solution. Conversely,
a solution supplies a full-rank matrix whose nonsingular
$r$-column submatrices are exactly $\B(M)$; this matrix represents
$M$. The encoding uses exact addition and multiplication tables
modulo five, so its satisfying assignments are precisely the
solutions of these field conditions. The harmless scaling
normalization and the arithmetic clauses are justified in
Appendix~\ref{sec:encoding}.
\end{proof}

For each final matroid, the SAT solver produces a DRAT refutation
of $\Phi_M$. DRAT-trim checks this refutation against the saved
formula \cite{drat}. By Proposition~\ref{prop:certificate}, a
successful check proves that $M$ is nonquinary. Together with
the 22 minor matrices, it certifies that $M$ is an excluded
minor. These certificates use the final basis records directly;
they do not need the deletion tables or their enumeration results.

At this point the logical roles are distinct. Complete parent
generation and exhaustive recovery show that every excluded
minor reaches the final list. Certificates show that every entry
of that list is an excluded minor. Exact isomorphism tests
determine how many different classes the list represents. We
now give the computational facts that establish these statements
for the present rank and size.

\section{The computation and the proof of the theorem}\label{sec:computation}

The completed run of 27 September 2026 records the outcomes in
Table~\ref{tab:stages}. Each row corresponds to a specified
mathematical set or operation in the preceding sections.

\begin{table}[htbp]
\centering
\begin{tabular}{p{22mm}p{95mm}r}
\toprule
Stage & Objects counted & Number\\
\midrule
Generation & Simple rank-four quinary ten-element parent classes
             & 58,756\\
Step 1 & Connected parent classes in $\mathcal P$ & 58,672\\
Step 2 & Labelled extensions, counted as pairs
         $(N,M)$ with $M\in\mathcal D(N)$ & 659\\
Step 3 & Classes after all contraction tests and isomorphism
         deduplication & 35\\
\bottomrule
\end{tabular}
\caption{The successive sets in the exhaustive classification.}
\label{tab:stages}
\end{table}

\subsection{Evidence that the construction is complete}

The represented rank-four generation retains 64,508
field-isomorphism classes at size ten. Abstract deduplication
gives the 58,756 simple parent classes in the first row. Exactly
84 are disconnected and are discarded by
Lemma~\ref{lem:parents}; the remaining 58,672 form $\mathcal P$.
Aligning the represented classes and closing their automorphism
orbits yields 71,547 labelled projective classes over the full
ten-element parent catalogue. The complete nine-element labelled
sets supply the other extension tables. All generated levels
are listed in Appendix~\ref{sec:reproduction}.

Every one of the 58,672 connected-parent problems finishes by
exhaustive natural join, with no unresolved case. The merged
logs account for every parent index exactly once, with each
parent either excluded for being disconnected or having its
finite problem exhausted. The shards have matching input and
critical-source fingerprints. By Propositions~\ref{prop:tables}
and~\ref{prop:local}, this yields exactly the pairs $(N,M)$
with $M\in\mathcal D(N)$: there are 659. This is a count of
parent-labelled extensions; it is not yet a count of isomorphism
classes of excluded minors.

Testing all contractions using the complete rank-three
catalogues, then deduplicating by isomorphism, leaves 35
classes. The independent incidence-graph check confirms that
these 35 are pairwise nonisomorphic. For each of them the
separate determinant formula has a checked DRAT refutation,
and all 22 saved minor matrices pass both modular Gaussian and
Sage verification. Thus there are 35 checked refutations and
$35\cdot22=770$ checked matrices.

\begin{proof}[Proof of Theorem~\ref{thm:main}]
To show that there are at most 35 classes, take any rank-four,
eleven-element quinary excluded minor $M$. By
Lemma~\ref{lem:parents}, a deletion of $M$ is isomorphic to a
parent $N\in\mathcal P$. The represented generation is complete
by Proposition~\ref{prop:generation}. After aligning the labels,
the recovery theorem, Proposition~\ref{prop:local}, places $M$
among the solutions of the search for $N$. Its tables are
complete by Proposition~\ref{prop:tables}, and the recorded
search exhausts them. Since every contraction of $M$ is quinary,
the complete contraction test of Lemma~\ref{lem:contractions}
retains it. Hence its class is among the 35 final classes.

To show that there are at least 35 classes, take the 35 final
representatives. The checked refutation for each proves its
nonrepresentability by Proposition~\ref{prop:certificate}.
Its 22 checked matrices prove representability of every
immediate minor, so it is an excluded minor. The independent
isomorphism check proves that the representatives belong to
35 different classes. Combining the two bounds proves the
classification.
\end{proof}

\subsection{Checks on the computation and consequences}

The same procedure reproduces the published rank-four catalogues
at sizes seven through ten: 5, 92, 219, and 130 classes,
respectively \cite{brettell}. Each comparison is by exact
matroid isomorphism with the published archive. At sizes seven
and eight the deletion-compatibility inequality $n>2r$ does not
hold, so those calibration runs impose basis exchange explicitly.
The rank-four nine-element target does satisfy the inequality.

A separate SAT enumeration of the finite table conditions agrees
with the join on every connected-parent problem for the
nine-element target. At the ten-element target, the complete
sets of 5,991 labelled deletion-minimal candidates agree. These
checks test the implementation; the proof of coverage is the
parent reduction and recovery argument above. The two enumeration
methods share their table construction, and the eleven-element
target is exhausted by the join method.

Duality gives a further classification immediately.
\begin{corollary}\label{cor:dual}
There are exactly 35 isomorphism classes of rank-seven,
eleven-element excluded minors for representability over $\F$.
\end{corollary}
\begin{proof}
Field representability is closed under duality, and duality
interchanges deletion and contraction. It therefore maps
excluded minors to excluded minors and gives a bijection between
the rank-four and rank-seven classes on eleven elements.
\end{proof}

Every matroid in the rank-four catalogue has a triangle, so none
is paving. This is a property of the resulting classification;
no paving restriction is used in the search. Basis and triangle
counts accompany the exact descriptions in
Appendix~\ref{sec:catalogue}.

\appendix
\section{Details of the representation and SAT encodings}\label{sec:encoding}

The main argument uses matrix normalization to enumerate complete
labelled representation sets, and an exact field encoding to certify
nonrepresentability. This appendix supplies the details of those
two operations. The last subsection describes the alternative SAT
enumeration used to cross-check the finite joins.

\subsection{A normal form for labelled representations}\label{sec:normalization}

Fix an ordered basis $B=(b_1,\ldots,b_r)$ of a labelled matroid
$Q$. Reduce a representation on $B$ so that its basis columns
form $I_r$, and call the remaining block $Z$. The support of
$Z$ is determined by fundamental exchanges:
\[
 Z_{ij}\ne0\quad\Longleftrightarrow\quad
       (B\setminus\{b_i\})\cup\{j\}\in\B(Q).
\]
Its support graph has one vertex for each row, one for each
nonbasis column, and an edge for each nonzero entry. Fix a
deterministic spanning forest of this graph. By row and column
scaling we make every forest entry equal to 1, compensating
on basis columns to preserve $I_r$.

\begin{lemma}[Normal form]\label{lem:normalform}
With the basis and support forest fixed, this normalization gives
one matrix for each labelled projective equivalence class.
\end{lemma}
\begin{proof}
To establish existence, root each nontrivial tree at a row vertex
and choose its scale. Traverse the tree: whenever the scale at
one endpoint of an edge is known, the nonzero entry on that edge
uniquely determines the scale at the other endpoint making the
entry 1. There is no cycle in the tree, so no conflicting choice
is imposed. Isolated vertices require no condition. Inverse
scaling of the basis columns restores the identity.

For uniqueness, two projectively equivalent representations
reduced on the same labelled basis differ by row and column
scalings: the row map must be diagonal because it fixes each
basis-column direction. Suppose both have forest entries 1.
On every tree edge, the row multiplier times the column
multiplier must be 1. Connectivity then forces a common row
multiplier throughout a component and its inverse on all column
vertices in that component. Their product is also 1 on every
nonforest edge, so all entries agree. This proves uniqueness,
including when the support graph is disconnected.
\end{proof}

In the labelled-orbit calculation of
Lemma~\ref{lem:labelled}, every new matrix is normalized before
comparison with those already seen. The lemma proves that this
comparison removes exactly projectively equivalent repetitions.
The full automorphism group of $Q$ is obtained from the colored
element--basis/nonbasis incidence graph described in
Section~\ref{sec:finish}. Its action on element vertices is
exactly $\Aut(Q)$, which supplies the generators used in the
orbit calculation.

\subsection{Compiling the determinant conditions to a formula}

For the formula $\Phi_M$ in Proposition~\ref{prop:certificate},
choose a basis and form $[I_r\mid Z]$ as above. The support is
fixed by the input matroid. Forest entries may be set to 1 by
Lemma~\ref{lem:normalform}, so only the other nonzero entries
need variables, each with domain $\{1,2,3,4\}$.

Let $S$ be an $r$-subset of the ground set and put
\[
 I(S)=\{i:b_i\notin S\},\qquad J(S)=S\setminus B.
\]
These sets have the same size. Expansion on the identity columns
shows that the full determinant on $S$ is, up to sign,
$\det Z[I(S),J(S)]$, with the empty determinant equal to 1.
It is therefore enough to impose
\[
 \det Z[I(S),J(S)]\ne0
 \quad\Longleftrightarrow\quad S\in\B(M)
 \qquad\text{for every }S\in\binom{E(M)}r.
\]
Every maximal minor is constrained, so these conditions determine
the full labelled matroid.

Here is an explicit finite propositional encoding of the field
arithmetic. For an expression $u$ with allowed values
$D\subseteq\F$, introduce indicators $p_{u,a}$ for $a\in D$,
requiring exactly one to be true. If $w=u+v$, add the clauses
\[
 \neg p_{u,a}\ \lor\ \neg p_{v,b}\ \lor\ p_{w,a+b}
\]
for every allowed pair $a,b$, where the sum is taken modulo five.
Use the corresponding clauses with $ab$ for multiplication.
The output indicators also satisfy the exactly-one condition.
Thus selected input values force precisely the correct output
value. Constants are substituted directly. A cached Laplace
expansion evaluates the determinants, and zero or nonzero
requirements constrain their indicators. Induction over the
arithmetic expressions proves that the formula represents
exactly the field conditions used in
Proposition~\ref{prop:certificate}.

For each listed matroid the saved evidence consists of its basis
record, the chosen basis and normalization, the DIMACS formula,
a DRAT refutation, and the proof-checker's output. The successful
refutation certifies the supplied formula's unsatisfiability;
the equivalence with representability is supplied by the
mathematical encoding above.

\subsection{The alternative SAT enumeration of tables}\label{sec:sat}

The natural-join search in Section~\ref{sec:tables} has an
alternative implementation as a propositional satisfiability
problem. It uses the same hyperplane coordinates and finite
relations $A_f$. A reduced binary decision diagram expresses
membership in each relation. If a node asks about $y$ and has
low and high children $l,h$, its variable $v$ is constrained by
\[
 v\ \Longleftrightarrow\ ((\neg y\land l)\lor(y\land h)).
\]
An exact four-clause encoding is
\[
 \begin{gathered}
 \neg v\lor y\lor l,\qquad \neg v\lor\neg y\lor h,\\
 v\lor y\lor\neg l,\qquad v\lor\neg y\lor\neg h.
 \end{gathered}
\]
Requiring each root enforces all table memberships. A blocking
clause excludes each represented full row in $A$. After an
assignment to the hyperplane coordinates is found, another
blocking clause excludes that assignment from subsequent
enumeration. A final UNSAT answer then establishes exhaustion
of this parent's solution set. A timeout is recorded as
\texttt{UNKNOWN}. The SAT/table-join comparisons reported in
Section~\ref{sec:computation} compare the complete assignments
obtained by these two enumeration procedures.

\section{Exact catalogue}\label{sec:catalogue}

All representatives have ground set $E=\{0,1,\ldots,10\}$ and
rank four. Each $M_i$ is specified by the listed family
$\HH_i$ of hyperplanes of size at least four, according to the
rule
\[
 \B(M_i)=\{S\in\binom E4:
                S\not\subseteq H\text{ for every }H\in\HH_i\}.
\]
Thus a four-set is a nonbasis exactly when it is contained in
a listed hyperplane. This specifies the whole basis family:
every dependent four-set of a rank-four matroid lies in a
hyperplane, while hyperplanes of size three contain no four-set
and may be omitted. The descriptions have been checked by
reconstructing the basis records from them.

The notation $M_i$ abbreviates the stable identifier
\texttt{GF5-r4-n11-} followed by $i$ in four digits. For example,
$M_1$ is \texttt{GF5-r4-n11-0001}. Proof directories are
zero-indexed, so its certificates are in
\texttt{matroid-0000}. Basis and triangle counts identify the
corresponding records in the machine-readable catalogue. Across
the catalogue there are between 210 and 262 bases, and the
triangle distribution is
\[
\begin{array}{c|rrrrrrrr}
\text{triangles} & 4 & 5 & 6 & 7 & 8 & 10 & 12 & 13\\ \hline
\text{classes} & 3 & 5 & 5 & 10 & 9 & 1 & 1 & 1
\end{array}

\]

\begingroup
\raggedright
% Generated by scripts/build_writeup.py from the checked catalogue.

\paragraph{$M_{1}$: 250 bases; 6 triangles.}
\(\{0,1,2,5,7,9\}\), \(\{0,1,4,10\}\), \(\{0,2,3,4,8\}\), \(\{0,2,6,10\}\), \(\{0,3,5,10\}\), \(\{0,3,6,9\}\), \(\{0,4,5,6\}\), \(\{0,6,7,8\}\), \(\{0,8,9,10\}\), \(\{1,2,3,4,5,6\}\), \(\{1,2,5,8\}\), \(\{1,2,5,10\}\), \(\{1,3,7,9\}\), \(\{1,4,7,9\}\), \(\{1,6,7,8,9,10\}\), \(\{2,3,4,7,10\}\), \(\{2,3,4,9\}\), \(\{2,6,7,8\}\), \(\{2,8,9,10\}\), \(\{3,6,7,8\}\), \(\{3,8,9,10\}\), \(\{4,5,6,7,8\}\), \(\{4,5,6,9\}\), \(\{4,5,6,10\}\), \(\{4,8,9,10\}\), \(\{5,8,9,10\}\).

\paragraph{$M_{2}$: 233 bases; 7 triangles.}
\(\{0,1,2,6\}\), \(\{0,1,4,7\}\), \(\{0,1,8,9,10\}\), \(\{0,2,3,4,5\}\), \(\{0,2,7,9\}\), \(\{0,3,6,9\}\), \(\{0,3,7,10\}\), \(\{0,4,6,8\}\), \(\{0,5,6,10\}\), \(\{0,5,7,8\}\), \(\{1,2,3,4,5,6,7\}\), \(\{1,2,6,8\}\), \(\{1,2,6,9\}\), \(\{1,2,6,10\}\), \(\{1,4,7,8\}\), \(\{1,4,7,9\}\), \(\{1,4,7,10\}\), \(\{2,3,4,5,8,9,10\}\), \(\{6,7,8,9,10\}\).

\paragraph{$M_{3}$: 252 bases; 7 triangles.}
\(\{0,1,2,5,7\}\), \(\{0,1,3,6\}\), \(\{0,1,4,9,10\}\), \(\{0,2,3,4,8\}\), \(\{0,2,6,9\}\), \(\{0,3,5,7\}\), \(\{0,4,5,6,7\}\), \(\{0,5,7,8,9\}\), \(\{0,5,7,10\}\), \(\{0,6,8,10\}\), \(\{1,2,3,4,5,6\}\), \(\{1,2,5,8\}\), \(\{1,2,5,9\}\), \(\{1,2,5,10\}\), \(\{1,3,6,7\}\), \(\{1,3,6,8,10\}\), \(\{1,3,6,9\}\), \(\{1,7,8,9\}\), \(\{2,3,4,7,10\}\), \(\{2,3,4,9\}\), \(\{2,6,8,10\}\), \(\{2,7,8,9\}\), \(\{3,5,9,10\}\), \(\{3,7,8,9\}\), \(\{4,5,6,8,10\}\), \(\{4,5,6,9\}\), \(\{4,7,8,9\}\), \(\{6,7,8,9,10\}\).

\paragraph{$M_{4}$: 250 bases; 5 triangles.}
\(\{0,1,2,5,7\}\), \(\{0,1,3,10\}\), \(\{0,1,4,8\}\), \(\{0,2,3,4,8,9\}\), \(\{0,2,6,10\}\), \(\{0,3,5,7\}\), \(\{0,4,5,7,8\}\), \(\{0,4,6,8\}\), \(\{0,4,8,10\}\), \(\{0,5,6,7,9\}\), \(\{0,5,7,10\}\), \(\{1,2,3,4,5,6\}\), \(\{1,2,5,8\}\), \(\{1,2,5,9\}\), \(\{1,2,5,10\}\), \(\{1,6,7,8,9,10\}\), \(\{2,3,4,7,10\}\), \(\{2,6,7,9\}\), \(\{3,5,8,10\}\), \(\{3,6,7,9\}\), \(\{4,5,9,10\}\), \(\{4,6,7,9\}\).

\paragraph{$M_{5}$: 247 bases; 7 triangles.}
\(\{0,1,2,6,8\}\), \(\{0,1,7,9\}\), \(\{0,2,3,4,5,7\}\), \(\{0,3,6,8,9,10\}\), \(\{0,4,6,8\}\), \(\{0,5,6,8\}\), \(\{0,6,7,8\}\), \(\{1,2,3,4,5,6\}\), \(\{1,2,6,7\}\), \(\{1,2,6,9\}\), \(\{1,2,6,10\}\), \(\{1,3,7,8\}\), \(\{1,4,7,10\}\), \(\{1,5,8,9,10\}\), \(\{2,3,4,5,8\}\), \(\{2,3,4,5,9\}\), \(\{2,3,4,5,10\}\), \(\{2,7,8,9,10\}\), \(\{4,6,7,9\}\), \(\{4,8,9,10\}\), \(\{5,6,7,10\}\).

\paragraph{$M_{6}$: 241 bases; 8 triangles.}
\(\{0,1,2,7,9,10\}\), \(\{0,1,6,8\}\), \(\{0,2,3,4,5,8\}\), \(\{0,3,7,9,10\}\), \(\{0,4,7,9,10\}\), \(\{0,5,6,7,9,10\}\), \(\{0,7,8,9,10\}\), \(\{1,2,3,4,5,6\}\), \(\{1,3,8,9\}\), \(\{1,4,8,10\}\), \(\{1,5,7,8\}\), \(\{2,3,4,5,7\}\), \(\{2,3,4,5,9\}\), \(\{2,3,4,5,10\}\), \(\{2,6,8,10\}\), \(\{3,6,7,8\}\), \(\{4,6,8,9\}\).

\paragraph{$M_{7}$: 249 bases; 7 triangles.}
\(\{0,1,2,6,8\}\), \(\{0,1,7,9,10\}\), \(\{0,2,3,4,5,7\}\), \(\{0,3,6,8,9\}\), \(\{0,4,6,8\}\), \(\{0,5,6,8,10\}\), \(\{0,6,7,8\}\), \(\{1,2,3,4,5,6\}\), \(\{1,2,6,7\}\), \(\{1,2,6,9\}\), \(\{1,2,6,10\}\), \(\{1,3,7,8\}\), \(\{1,4,8,10\}\), \(\{1,5,8,9\}\), \(\{2,3,4,5,8\}\), \(\{2,3,4,5,9\}\), \(\{2,3,4,5,10\}\), \(\{2,7,8,9,10\}\), \(\{3,7,9,10\}\), \(\{4,6,7,9,10\}\), \(\{5,7,9,10\}\).

\paragraph{$M_{8}$: 252 bases; 7 triangles.}
\(\{0,1,2,5,7\}\), \(\{0,1,3,6\}\), \(\{0,1,8,9\}\), \(\{0,2,3,4,8\}\), \(\{0,3,5,7\}\), \(\{0,4,5,6,7\}\), \(\{0,5,7,8\}\), \(\{0,5,7,9,10\}\), \(\{0,6,8,10\}\), \(\{1,2,3,4,5,6\}\), \(\{1,2,5,8,9\}\), \(\{1,2,5,10\}\), \(\{1,3,6,7\}\), \(\{1,3,6,8,9\}\), \(\{1,3,6,10\}\), \(\{1,4,8,9\}\), \(\{1,7,8,9,10\}\), \(\{2,3,4,7,9,10\}\), \(\{4,5,6,8\}\), \(\{4,5,6,9\}\), \(\{4,5,6,10\}\), \(\{6,7,9,10\}\).

\paragraph{$M_{9}$: 255 bases; 6 triangles.}
\(\{0,1,2,5,7\}\), \(\{0,1,3,8\}\), \(\{0,1,6,9,10\}\), \(\{0,2,3,4,8\}\), \(\{0,2,6,9\}\), \(\{0,3,5,8\}\), \(\{0,3,6,8,9\}\), \(\{0,3,7,8,10\}\), \(\{0,4,5,6,9\}\), \(\{0,6,7,9\}\), \(\{1,2,3,4,5,6\}\), \(\{1,2,5,8,9\}\), \(\{1,2,5,10\}\), \(\{1,3,7,9\}\), \(\{1,7,8,10\}\), \(\{2,3,4,7\}\), \(\{2,3,4,9\}\), \(\{2,3,4,10\}\), \(\{2,6,7,8,10\}\), \(\{3,5,9,10\}\), \(\{4,5,6,7\}\), \(\{4,5,6,8\}\), \(\{4,5,6,10\}\), \(\{4,7,8,9,10\}\), \(\{5,7,8,10\}\).

\paragraph{$M_{10}$: 239 bases; 8 triangles.}
\(\{0,1,2,6,8\}\), \(\{0,1,4,9,10\}\), \(\{0,1,5,7\}\), \(\{0,2,3,4,5\}\), \(\{0,2,7,9\}\), \(\{0,3,6,8,9\}\), \(\{0,3,7,10\}\), \(\{0,4,6,7,8\}\), \(\{0,5,6,8\}\), \(\{0,6,8,10\}\), \(\{1,2,3,4,5,6,7\}\), \(\{1,2,6,9\}\), \(\{1,2,6,10\}\), \(\{1,5,7,8\}\), \(\{1,5,7,9\}\), \(\{1,5,7,10\}\), \(\{2,3,4,5,8,10\}\), \(\{2,3,4,5,9\}\), \(\{4,6,7,9\}\), \(\{4,6,7,10\}\), \(\{5,6,9,10\}\), \(\{7,8,9,10\}\).

\paragraph{$M_{11}$: 212 bases; 10 triangles.}
\(\{0,1,2,6,8,9\}\), \(\{0,1,4,10\}\), \(\{0,1,5,7\}\), \(\{0,2,3,4,5\}\), \(\{0,2,7,10\}\), \(\{0,3,6,10\}\), \(\{0,3,7,9\}\), \(\{0,4,6,7\}\), \(\{0,5,8,10\}\), \(\{1,2,3,4,5,6,7,8\}\), \(\{1,2,6,8,10\}\), \(\{1,5,7,9\}\), \(\{1,5,7,10\}\), \(\{2,3,4,5,9,10\}\), \(\{4,6,7,9\}\), \(\{4,6,7,10\}\), \(\{7,8,9,10\}\).

\paragraph{$M_{12}$: 242 bases; 8 triangles.}
\(\{0,1,2,6,8\}\), \(\{0,1,3,10\}\), \(\{0,1,4,9\}\), \(\{0,1,5,7\}\), \(\{0,2,3,4,5,9\}\), \(\{0,4,6,7,9\}\), \(\{0,4,8,9\}\), \(\{0,4,9,10\}\), \(\{0,5,8,10\}\), \(\{1,2,3,4,5,6,7\}\), \(\{1,2,6,9,10\}\), \(\{1,5,7,8\}\), \(\{1,5,7,9\}\), \(\{1,5,7,10\}\), \(\{2,3,4,5,8\}\), \(\{2,3,4,5,10\}\), \(\{3,7,8,9,10\}\), \(\{4,6,7,8\}\), \(\{4,6,7,10\}\).

\paragraph{$M_{13}$: 254 bases; 5 triangles.}
\(\{0,1,2,5,7\}\), \(\{0,1,4,8\}\), \(\{0,2,3,4,8,9\}\), \(\{0,3,5,7\}\), \(\{0,4,5,7,8\}\), \(\{0,4,6,8\}\), \(\{0,4,8,10\}\), \(\{0,5,6,7,9,10\}\), \(\{1,2,3,4,5,6\}\), \(\{1,2,5,8\}\), \(\{1,2,5,9\}\), \(\{1,2,5,10\}\), \(\{1,3,7,8\}\), \(\{1,6,8,9,10\}\), \(\{2,3,4,7,10\}\), \(\{2,6,9,10\}\), \(\{3,6,9,10\}\), \(\{4,6,9,10\}\).

\paragraph{$M_{14}$: 257 bases; 4 triangles.}
\(\{0,1,2,5,7,9\}\), \(\{0,1,4,10\}\), \(\{0,2,3,4,8\}\), \(\{0,3,6,7,9,10\}\), \(\{0,4,7,9\}\), \(\{0,5,6,8\}\), \(\{0,7,8,9\}\), \(\{1,2,3,4,5,6\}\), \(\{1,2,5,8\}\), \(\{1,2,5,10\}\), \(\{1,3,7,8\}\), \(\{1,6,8,9,10\}\), \(\{2,3,4,7\}\), \(\{2,3,4,9\}\), \(\{2,3,4,10\}\), \(\{2,6,9,10\}\), \(\{4,5,7,8,10\}\), \(\{4,6,9,10\}\), \(\{5,6,9,10\}\).

\paragraph{$M_{15}$: 262 bases; 4 triangles.}
\(\{0,1,2,5,7\}\), \(\{0,1,4,10\}\), \(\{0,2,3,4,8\}\), \(\{0,3,5,7,10\}\), \(\{0,4,5,7\}\), \(\{0,5,6,7,8,9\}\), \(\{1,2,3,4,5,6\}\), \(\{1,2,5,8\}\), \(\{1,2,5,9\}\), \(\{1,2,5,10\}\), \(\{1,3,7,9\}\), \(\{1,4,7,8\}\), \(\{1,6,8,9,10\}\), \(\{2,3,4,7\}\), \(\{2,3,4,9,10\}\), \(\{2,6,8,9\}\), \(\{2,7,8,10\}\), \(\{3,6,8,9\}\), \(\{4,6,7,10\}\), \(\{4,6,8,9\}\).

\paragraph{$M_{16}$: 210 bases; 13 triangles.}
\(\{0,1,2,7\}\), \(\{0,1,4,9,10\}\), \(\{0,1,5,8\}\), \(\{0,2,3,4,5,6\}\), \(\{0,2,8,9\}\), \(\{0,3,7,9\}\), \(\{0,3,8,10\}\), \(\{0,4,7,8\}\), \(\{0,6,7,10\}\), \(\{1,2,3,4,5,6,7,8\}\), \(\{1,2,7,9\}\), \(\{1,2,7,10\}\), \(\{1,5,8,9\}\), \(\{1,5,8,10\}\), \(\{2,3,4,5,6,9\}\), \(\{2,3,4,5,6,10\}\), \(\{4,7,8,9\}\), \(\{4,7,8,10\}\), \(\{5,7,9,10\}\), \(\{6,8,9,10\}\).

\paragraph{$M_{17}$: 245 bases; 7 triangles.}
\(\{0,1,2,5,7,9\}\), \(\{0,1,4,8\}\), \(\{0,2,3,4,8,9\}\), \(\{0,2,6,9\}\), \(\{0,2,9,10\}\), \(\{0,3,5,7\}\), \(\{0,4,5,7,8\}\), \(\{0,4,6,8\}\), \(\{0,4,8,10\}\), \(\{0,5,6,7,10\}\), \(\{1,2,3,4,5,6\}\), \(\{1,2,5,8,10\}\), \(\{1,3,7,8,9\}\), \(\{1,4,7,9\}\), \(\{1,6,7,9\}\), \(\{1,7,9,10\}\), \(\{2,3,4,7,10\}\), \(\{3,5,8,9\}\), \(\{3,6,8,9,10\}\).

\paragraph{$M_{18}$: 242 bases; 6 triangles.}
\(\{0,1,2,5,8,10\}\), \(\{0,1,3,10\}\), \(\{0,1,4,9,10\}\), \(\{0,1,6,10\}\), \(\{0,1,7,10\}\), \(\{0,2,3,4,9\}\), \(\{0,3,5,7\}\), \(\{0,4,5,6,9\}\), \(\{0,4,7,9\}\), \(\{0,4,8,9\}\), \(\{0,6,7,8\}\), \(\{1,2,3,4,5,6,7\}\), \(\{1,2,5,9\}\), \(\{1,7,8,9\}\), \(\{2,3,4,8\}\), \(\{2,3,4,10\}\), \(\{2,7,9,10\}\), \(\{3,5,7,8\}\), \(\{3,5,7,9\}\), \(\{3,5,7,10\}\), \(\{3,6,8,9,10\}\), \(\{4,5,6,8\}\), \(\{4,5,6,10\}\), \(\{4,7,8,10\}\).

\paragraph{$M_{19}$: 244 bases; 8 triangles.}
\(\{0,1,2,6,7,9\}\), \(\{0,1,3,9\}\), \(\{0,1,4,8,9\}\), \(\{0,1,5,9,10\}\), \(\{0,2,3,4,5,7\}\), \(\{0,2,7,8\}\), \(\{0,2,7,10\}\), \(\{0,6,8,10\}\), \(\{1,2,3,4,5,6\}\), \(\{1,2,6,8\}\), \(\{1,2,6,10\}\), \(\{1,3,8,10\}\), \(\{2,3,4,5,8\}\), \(\{2,3,4,5,9\}\), \(\{2,3,4,5,10\}\), \(\{2,8,9,10\}\), \(\{3,6,7,8,9\}\), \(\{4,6,7,9,10\}\), \(\{5,6,7,9\}\), \(\{5,7,8,10\}\).

\paragraph{$M_{20}$: 256 bases; 6 triangles.}
\(\{0,1,2,5,7\}\), \(\{0,1,3,9,10\}\), \(\{0,1,4,6\}\), \(\{0,2,3,4,6\}\), \(\{0,3,5,7\}\), \(\{0,4,5,6,7\}\), \(\{0,4,6,8,10\}\), \(\{0,4,6,9\}\), \(\{0,5,7,8,9\}\), \(\{0,5,7,10\}\), \(\{1,2,3,4,5\}\), \(\{1,2,5,6\}\), \(\{1,2,5,8,9\}\), \(\{1,2,5,10\}\), \(\{1,3,6,7,8\}\), \(\{1,4,7,9\}\), \(\{1,4,8,10\}\), \(\{2,3,4,7\}\), \(\{2,3,4,8,10\}\), \(\{2,3,4,9\}\), \(\{2,6,7,9,10\}\), \(\{3,5,6,10\}\), \(\{3,5,8,9\}\), \(\{4,5,8,9,10\}\), \(\{4,7,8,10\}\), \(\{5,6,8,9\}\).

\paragraph{$M_{21}$: 226 bases; 8 triangles.}
\(\{0,1,2,6,8,9\}\), \(\{0,1,4,7\}\), \(\{0,1,5,10\}\), \(\{0,2,3,4,5\}\), \(\{0,2,7,10\}\), \(\{0,3,7,9\}\), \(\{0,3,8,10\}\), \(\{0,4,6,10\}\), \(\{0,5,7,8\}\), \(\{1,2,3,4,5,6,7\}\), \(\{1,2,6,10\}\), \(\{1,3,8,10\}\), \(\{1,4,7,8\}\), \(\{1,4,7,9\}\), \(\{1,4,7,10\}\), \(\{2,3,4,5,8,9,10\}\), \(\{2,7,8,9\}\), \(\{3,6,8,10\}\), \(\{3,7,8,10\}\), \(\{6,7,9,10\}\).

\paragraph{$M_{22}$: 261 bases; 4 triangles.}
\(\{0,1,2,5,7\}\), \(\{0,1,6,8,9\}\), \(\{0,2,3,4,8\}\), \(\{0,2,6,9\}\), \(\{0,3,5,7\}\), \(\{0,3,6,9\}\), \(\{0,4,5,7\}\), \(\{0,4,6,9\}\), \(\{0,5,6,7,9,10\}\), \(\{0,5,7,8\}\), \(\{1,2,3,4,5,6\}\), \(\{1,2,5,8,10\}\), \(\{1,2,5,9\}\), \(\{1,3,7,10\}\), \(\{1,4,7,9\}\), \(\{2,3,4,7\}\), \(\{2,3,4,9,10\}\), \(\{2,7,8,9\}\), \(\{3,5,8,9\}\), \(\{3,6,7,8\}\), \(\{4,7,8,10\}\).

\paragraph{$M_{23}$: 220 bases; 8 triangles.}
\(\{0,1,2,6\}\), \(\{0,1,4,7\}\), \(\{0,1,5,9\}\), \(\{0,1,8,10\}\), \(\{0,2,3,4,5,10\}\), \(\{0,2,8,9\}\), \(\{0,3,6,8\}\), \(\{0,5,7,8\}\), \(\{0,6,7,9,10\}\), \(\{1,2,3,4,5,6,7,8\}\), \(\{1,2,6,9\}\), \(\{1,2,6,10\}\), \(\{1,3,9,10\}\), \(\{1,4,7,9\}\), \(\{1,4,7,10\}\), \(\{2,3,4,5,9\}\), \(\{3,6,8,9\}\), \(\{3,6,8,10\}\), \(\{4,8,9,10\}\), \(\{5,7,8,9\}\), \(\{5,7,8,10\}\).

\paragraph{$M_{24}$: 220 bases; 8 triangles.}
\(\{0,1,2,6\}\), \(\{0,1,4,7\}\), \(\{0,1,8,9,10\}\), \(\{0,2,3,4,5\}\), \(\{0,2,7,9\}\), \(\{0,3,6,9\}\), \(\{0,3,7,10\}\), \(\{0,4,6,8\}\), \(\{0,5,6,10\}\), \(\{0,5,7,8\}\), \(\{1,2,3,4,5,6,7,8\}\), \(\{1,2,6,9\}\), \(\{1,2,6,10\}\), \(\{1,4,7,9\}\), \(\{1,4,7,10\}\), \(\{2,3,4,5,9,10\}\), \(\{4,6,8,9\}\), \(\{4,6,8,10\}\), \(\{5,7,8,9\}\), \(\{5,7,8,10\}\), \(\{6,7,9,10\}\).

\paragraph{$M_{25}$: 243 bases; 7 triangles.}
\(\{0,1,2,6\}\), \(\{0,1,4,7\}\), \(\{0,1,8,9,10\}\), \(\{0,2,3,4,5\}\), \(\{0,2,7,10\}\), \(\{0,2,8,9\}\), \(\{0,3,6,10\}\), \(\{0,3,8,9\}\), \(\{0,4,6,8,9\}\), \(\{0,5,7,8,9\}\), \(\{1,2,3,4,5,6,7\}\), \(\{1,2,6,8\}\), \(\{1,2,6,9\}\), \(\{1,2,6,10\}\), \(\{1,4,7,8\}\), \(\{1,4,7,9\}\), \(\{1,4,7,10\}\), \(\{2,3,4,5,8\}\), \(\{2,3,4,5,9,10\}\), \(\{3,7,8,10\}\), \(\{5,6,8,10\}\), \(\{6,7,9,10\}\).

\paragraph{$M_{26}$: 243 bases; 7 triangles.}
\(\{0,1,2,6\}\), \(\{0,1,4,7\}\), \(\{0,1,8,9,10\}\), \(\{0,2,3,4,5\}\), \(\{0,2,8,9\}\), \(\{0,3,6,8,9\}\), \(\{0,3,7,10\}\), \(\{0,4,6,10\}\), \(\{0,4,8,9\}\), \(\{0,5,7,8,9\}\), \(\{1,2,3,4,5,6,7\}\), \(\{1,2,6,8\}\), \(\{1,2,6,9\}\), \(\{1,2,6,10\}\), \(\{1,4,7,8\}\), \(\{1,4,7,9\}\), \(\{1,4,7,10\}\), \(\{2,3,4,5,8\}\), \(\{2,3,4,5,9,10\}\), \(\{2,7,8,10\}\), \(\{5,6,8,10\}\), \(\{6,7,9,10\}\).

\paragraph{$M_{27}$: 228 bases; 12 triangles.}
\(\{0,1,2,7\}\), \(\{0,1,4,9\}\), \(\{0,1,5,8,10\}\), \(\{0,2,3,4,5,6\}\), \(\{0,2,8,9,10\}\), \(\{0,3,7,9\}\), \(\{0,3,8,10\}\), \(\{0,4,7,8,10\}\), \(\{0,6,8,10\}\), \(\{1,2,3,4,5,6,7\}\), \(\{1,2,7,8\}\), \(\{1,2,7,9\}\), \(\{1,2,7,10\}\), \(\{1,3,9,10\}\), \(\{1,6,8,9\}\), \(\{2,3,4,5,6,8\}\), \(\{2,3,4,5,6,9\}\), \(\{2,3,4,5,6,10\}\), \(\{5,7,8,9\}\), \(\{6,7,9,10\}\).

\paragraph{$M_{28}$: 248 bases; 5 triangles.}
\(\{0,1,2,5,8\}\), \(\{0,1,7,9\}\), \(\{0,2,3,4\}\), \(\{0,2,7,9\}\), \(\{0,3,6,7,9,10\}\), \(\{0,4,5,6\}\), \(\{0,4,7,8,9\}\), \(\{0,5,7,9\}\), \(\{1,2,3,4,5,6,7\}\), \(\{1,2,5,9,10\}\), \(\{1,3,8,10\}\), \(\{2,3,4,8\}\), \(\{2,3,4,9\}\), \(\{2,3,4,10\}\), \(\{2,6,8,9\}\), \(\{3,6,7,8\}\), \(\{4,5,6,8,10\}\), \(\{4,5,6,9\}\).

\paragraph{$M_{29}$: 216 bases; 8 triangles.}
\(\{0,1,2,9\}\), \(\{0,1,6,7,8,10\}\), \(\{0,2,3,4,5,10\}\), \(\{0,3,8,9\}\), \(\{0,4,7,9\}\), \(\{0,5,6,9\}\), \(\{1,2,3,4,5,6,7,8\}\), \(\{1,3,9,10\}\), \(\{1,6,7,8,9\}\), \(\{2,3,4,5,9\}\), \(\{2,7,9,10\}\), \(\{4,6,9,10\}\), \(\{5,8,9,10\}\).

\paragraph{$M_{30}$: 256 bases; 5 triangles.}
\(\{0,1,2,5,7\}\), \(\{0,1,4,9,10\}\), \(\{0,2,3,4,8\}\), \(\{0,2,6,9,10\}\), \(\{0,3,5,9,10\}\), \(\{0,4,5,6\}\), \(\{0,6,7,8\}\), \(\{0,7,9,10\}\), \(\{0,8,9,10\}\), \(\{1,2,3,4,5,6\}\), \(\{1,2,5,8,9\}\), \(\{1,2,5,10\}\), \(\{1,6,7,8,10\}\), \(\{2,3,4,7,10\}\), \(\{2,3,4,9\}\), \(\{2,6,7,8\}\), \(\{3,6,7,8,9\}\), \(\{4,5,6,7,8\}\), \(\{4,5,6,9\}\), \(\{4,5,6,10\}\).

\paragraph{$M_{31}$: 239 bases; 7 triangles.}
\(\{0,1,2,6,8,9\}\), \(\{0,1,4,7\}\), \(\{0,2,3,4,5,10\}\), \(\{0,3,6,8\}\), \(\{0,4,6,8\}\), \(\{0,5,6,8\}\), \(\{0,5,7,9\}\), \(\{0,6,7,8,10\}\), \(\{1,2,3,4,5,6,7\}\), \(\{1,2,6,10\}\), \(\{1,3,9,10\}\), \(\{1,4,7,8\}\), \(\{1,4,7,9\}\), \(\{1,4,7,10\}\), \(\{2,3,4,5,8\}\), \(\{2,3,4,5,9\}\), \(\{2,7,9,10\}\), \(\{3,7,8,9\}\), \(\{4,6,9,10\}\), \(\{5,8,9,10\}\).

\paragraph{$M_{32}$: 241 bases; 7 triangles.}
\(\{0,1,2,5,8\}\), \(\{0,1,3,7\}\), \(\{0,1,4,9,10\}\), \(\{0,2,3,4\}\), \(\{0,2,6,7,9\}\), \(\{0,3,5,8,9\}\), \(\{0,3,6,10\}\), \(\{0,4,5,6,8\}\), \(\{0,5,7,8\}\), \(\{0,5,8,10\}\), \(\{1,2,3,4,5,6,7\}\), \(\{1,2,5,9\}\), \(\{1,2,5,10\}\), \(\{1,3,7,8,9\}\), \(\{1,3,7,10\}\), \(\{2,3,4,8,9\}\), \(\{2,3,4,10\}\), \(\{2,6,7,8,10\}\), \(\{3,6,8,9\}\), \(\{3,8,9,10\}\), \(\{4,5,6,9\}\), \(\{4,5,6,10\}\), \(\{5,7,9,10\}\).

\paragraph{$M_{33}$: 241 bases; 6 triangles.}
\(\{0,1,2,5,8,9\}\), \(\{0,1,3,10\}\), \(\{0,2,3,4\}\), \(\{0,2,6,10\}\), \(\{0,3,6,7,9\}\), \(\{0,4,5,6\}\), \(\{0,4,7,8\}\), \(\{0,4,9,10\}\), \(\{0,5,7,10\}\), \(\{1,2,3,4,5,6,7\}\), \(\{1,2,5,10\}\), \(\{1,3,8,9\}\), \(\{1,4,8,9,10\}\), \(\{1,6,8,9\}\), \(\{1,7,8,9\}\), \(\{2,3,4,8\}\), \(\{2,3,4,9,10\}\), \(\{3,6,7,8,10\}\), \(\{4,5,6,8\}\), \(\{4,5,6,9,10\}\), \(\{4,7,9,10\}\).

\paragraph{$M_{34}$: 258 bases; 5 triangles.}
\(\{0,1,2,5,7,9\}\), \(\{0,1,4,10\}\), \(\{0,2,3,4,8\}\), \(\{0,3,5,7\}\), \(\{0,3,6,9,10\}\), \(\{0,4,5,6,7\}\), \(\{0,5,7,8\}\), \(\{0,5,7,10\}\), \(\{1,2,3,4,5,6\}\), \(\{1,2,5,8,10\}\), \(\{1,3,7,9\}\), \(\{1,4,7,9\}\), \(\{1,6,7,8,9\}\), \(\{1,7,9,10\}\), \(\{2,3,4,7\}\), \(\{2,3,4,9\}\), \(\{2,3,4,10\}\), \(\{2,6,7,10\}\), \(\{3,5,8,9\}\), \(\{3,7,8,10\}\), \(\{4,5,6,8\}\), \(\{4,5,6,9\}\), \(\{4,5,6,10\}\), \(\{4,8,9,10\}\).

\paragraph{$M_{35}$: 238 bases; 8 triangles.}
\(\{0,1,2,6,8\}\), \(\{0,1,5,7,9\}\), \(\{0,2,3,4,5,10\}\), \(\{0,3,6,8,9\}\), \(\{0,4,6,7,8\}\), \(\{0,5,6,8\}\), \(\{0,6,8,10\}\), \(\{1,2,3,4,5,6,7\}\), \(\{1,2,6,9\}\), \(\{1,2,6,10\}\), \(\{1,3,9,10\}\), \(\{1,5,7,8,10\}\), \(\{2,3,4,5,8\}\), \(\{2,3,4,5,9\}\), \(\{2,7,9,10\}\), \(\{4,6,7,9\}\), \(\{4,6,7,10\}\), \(\{4,8,9,10\}\), \(\{5,6,9,10\}\).

\par\endgroup

\section{Computational records and reproduction}\label{sec:reproduction}

This appendix identifies the saved evidence for each step of
the proof. All paths are relative to the repository root.
The exact catalogue is
\path{results/r4-n11/excluded.jsonl}, with SHA-256 digest
\begin{quote}\small\ttfamily
7940989ae051b32496dbfd7e4717d405431bcfd89da69d0dd522517552c76441.
\end{quote}
Its equivalent hyperplane descriptions are in
\path{results/r4-n11/hyperplanes.json}.

\subsection{Parent coverage and exhaustive search}

Table~\ref{tab:parents} gives the completed represented levels.
The rank-four size-four entry is the seed $I_4$; rank-three
generation starts from $I_3$. The counts distinguish represented
classes, which must be retained during generation, from abstract
matroid classes used to index the parent searches.

\begin{table}[htbp]
\centering
\begin{tabular}{rrrrr}
\toprule
Size & Rank 3 field & Rank 3 abstract & Rank 4 field & Rank 4 abstract \\
\midrule
4 & 2 & 2 & 1 & 1 \\
5 & 4 & 4 & 3 & 3 \\
6 & 11 & 9 & 12 & 11 \\
7 & 22 & 18 & 58 & 44 \\
8 & 42 & 34 & 460 & 303 \\
9 & 92 & 82 & 5,001 & 3,763 \\
10 & 174 & 168 & 64,508 & 58,756 \\
\bottomrule
\end{tabular}

\caption{Complete simple represented catalogues. ``Field'' counts
field-isomorphism classes; ``abstract'' counts matroid classes.}
\label{tab:parents}
\end{table}

The main scripts correspond to the proof as follows.
\begin{description}\raggedright
\item[Parent generation.] \path{scripts/generate_parents.py}
generates all represented rank-four and rank-three levels through
size ten. The completed abstract lists index the parents and
supply the contraction comparisons.
\item[Labelled representations.]
\path{scripts/prepare_representations.py} aligns the represented
classes and closes under automorphisms. The normalization and
point-extension tables are implemented in
\path{gf5minors/projective.py}.
\item[Recovery and contraction testing.]
\path{scripts/search_compatible_deletions.py} builds the tables
for each parent and tests contractions of the recovered
matroids. The exact join is in
\path{gf5minors/table_join.py}; the alternative SAT solver is in
\path{gf5minors/table_sat.py}.
\item[Merging and counting.]
\path{scripts/merge_search.py} checks parent coverage and removes
isomorphic repetitions. \path{scripts/final_audit.py} checks
certificate coverage and distinctness by canonical incidence
graphs.
\end{description}

The four ranges in Table~\ref{tab:shards} partition
$[0,58756)$. Their logs record either exhaustion or the
disconnected-parent exclusion for every index. The merge checks
the ranges for gaps and overlaps, terminal statuses, and
agreement of input and critical-source fingerprints.

\begin{table}[htbp]
\centering
\begin{tabular}{rrrrr}
\toprule
Start & Stop & Parents & Candidates & Seconds \\
\midrule
0 & 14,689 & 14,689 & 514 & 1794.54 \\
14,689 & 29,378 & 14,689 & 106 & 1629.19 \\
29,378 & 44,067 & 14,689 & 29 & 1430.24 \\
44,067 & 58,756 & 14,689 & 10 & 1299.00 \\
\bottomrule
\end{tabular}

\caption{Exhaustive target-search ranges. Candidate counts precede
contraction filtering and global isomorphism deduplication.}
\label{tab:shards}
\end{table}

The global coverage record is
\path{results/r4-n11/summary.json}. Detailed parent logs and
input manifests are in \path{runs/r4n11/shard-00/} through
\path{runs/r4n11/shard-03/}. The final audit and statistics are
\path{results/r4-n11/verification.json} and
\path{results/r4-n11/statistics.json}. The target search used
four Slurm tasks, each requesting one CPU and 4 GB of memory.
The longest range took about 30 minutes.

\subsection{Certificates and their replay}

\path{scripts/certify_nonrepresentation.py} constructs the
separate determinant formulas and checks their DRAT proofs.
The 35 proof directories under
\path{results/r4-n11/nonrepresentation/} are named
\texttt{matroid-0000} through \texttt{matroid-0034}. For $M_i$,
the directory index is $i-1$. Its files
\texttt{representation.cnf} and \texttt{proof.drat} can be
replayed with
\begin{verbatim}
vendor/drat-trim/drat-trim FORMULA PROOF
\end{verbatim}
A successful check prints \texttt{s VERIFIED}.

\path{scripts/certify_minors.py} constructs and checks the
immediate-minor matrices. All 770 are stored in
\path{results/r4-n11/minors/minor-witnesses.jsonl}.
The remaining column labels are in increasing parent order,
renumbered to $0,\ldots,9$. The pure modular checker
\texttt{gf5minors.matroid.verify\_representation} can check
these matrices directly from the basis records, independently
of the search and the represented catalogues.

The completed computation records Sage 10.9 and, in a separate
environment, Python 3.11.16, python-sat 1.9.dev15, and NumPy
2.4.6. Sage supplies represented extensions, isomorphism,
and graph operations \cite{sage}. The proof checker is pinned
to DRAT-trim commit
\texttt{2e3b2dc0ecf938addbd779d42877b6ed69d9a985}.
The computational proof relies on the correctness of parent
generation, representation orbits, table construction,
exhaustive enumeration, and isomorphism handling. The DRAT
checks certify the final determinant formulas; coverage of the
classification is supplied by the exhaustive construction.

On the original cluster, run computation through Slurm with
the existing \texttt{my-sage} and \texttt{matroid1}
environments. \path{slurm/research.slurm} selects the environment
and caps numerical-library threads at one. Full reproduction
commands, in dependency order, are given in
\path{docs/REPRODUCING.md}. Use fresh output directories for new
runs. The reproducibility archive
\path{results/r4-n11-reproducibility.tar.gz} and its adjacent
JSON manifest preserve the original result and its inputs;
that archive predates this exposition.

\begin{thebibliography}{9}

\bibitem{brettell}
N. Brettell,
\emph{The excluded minors for $\mathrm{GF}(5)$-representable matroids
on ten elements},
Advances in Applied Mathematics \textbf{166} (2025), 102864.
\href{https://doi.org/10.1016/j.aam.2025.102864}{doi:10.1016/j.aam.2025.102864}.
Preprint: \href{https://arxiv.org/abs/2307.14614}{arXiv:2307.14614},
version 2, 2024.

\bibitem{brettell-pendavingh}
N. Brettell and R. Pendavingh,
\emph{Computing excluded minors for classes of matroids representable
over partial fields},
Electronic Journal of Combinatorics \textbf{31}(3) (2024), P3.20.
\href{https://arxiv.org/abs/2302.13175}{arXiv:2302.13175}.

\bibitem{oxley}
J. Oxley,
\emph{Matroid Theory}, second edition,
Oxford Graduate Texts in Mathematics 21,
Oxford University Press, 2011.

\bibitem{sage}
The Sage Developers,
\emph{SageMath, the Sage Mathematics Software System},
version 10.9, used in the computation.
\href{https://doc.sagemath.org/html/en/reference/matroids/sage/matroids/linear_matroid.html}{Linear matroids reference manual}.

\bibitem{drat}
N. Wetzler, M. J. H. Heule, and W. A. Hunt, Jr.,
\emph{DRAT-trim: Efficient checking and trimming using expressive
clausal proofs},
Theory and Applications of Satisfiability Testing (SAT 2014),
Lecture Notes in Computer Science \textbf{8561}, 422--429.
\href{https://www.cs.utexas.edu/~marijn/drat-trim/}{Software and documentation}.

\end{thebibliography}
\end{document}
